%! Unit Launch: Study Design — Unit 1, Lesson 1.0 — Homework
% PILOT SHAPE (2026-09-06): modeled on AP Statistics lesson 1.4 minus its AP Practice page.
% This is the lesson's SCORED individual practice and the LAST component in the packet.
% Students START IT IN CLASS in the last ten minutes of the period, alone, while the
% teacher circulates for the second formative read; they finish it at home. Packet pages are
% the default; a Desmos activity may replace them on a given day (decided at Close & Assign).
% THIRD CONTEXT: the notes teach on the Riverbend basketball roster, the Guided Practice runs
% on the Student Life Survey, and these items are on the Northgate Recreation Center — recall
% is not the skill.
% TWO PAGES MAX (user decision 2026-09-06): two boxes — items 1-2 on page 1, items 3-6 and
% the spiralbox on page 2 — so no item breaks across a page.
% Item 2's Member ID and Age rows are the deliberate CONTRAST PAIR: two columns of digits, one
% a label and one an amount. Item 6 is the target misconception head-on (the mean Member ID).
% THE DATA — Northgate Recreation Center, 8 members: Member ID 104,118,127,133,145,152,166,179
%   -> 1124/8 = 140.5; visits 6,2,9,4,14,1,8,4 -> 48/8 = 6; group class 5/8 = 62.5%;
%   pool survey 45/300 = 15%; 0.15 x 2000 = 300.
% PAGE LOCKSTEP: the key is generated from this file; every work block is byte-identical.
\documentclass[12pt]{article}
\usepackage{saar-article}
\usepackage{saar-boxes}
\newcommand{\TallMath}[1]{$\displaystyle #1\rule[-1.4em]{0pt}{3.2em}$}
% Word bank strip for every fill-in cluster (user request 2026-09-06). Defined per document;
% the key inherits it and prints the same strip, so heights match.
\newcommand{\wordbank}[1]{\par\vspace{2pt}\noindent\fcolorbox{goldacc}{goldbg}{\parbox{\dimexpr\linewidth-2\fboxsep-2\fboxrule\relax}{\small\textbf{Word bank:} #1}}\par\vspace{4pt}}

\begin{document}

\pageheader{Unit 1, Lesson 1.0}{Homework}
% NAMESTRIP: no \namedateperiod here. The student writes their name once, on the
% cover this component is stapled behind.

\vspace{0.12in}

% Authored identically in homework/ and homework_key/.
\begin{remindbox}
\small
\textbf{This is your graded homework.} It is scored out of 12 and due at the start of the
first class after two study halls. You will start it in the last minutes of the period, so ask about anything that stalls
you before you leave, then finish it on your own.
\end{remindbox}

\vspace{0.08in}

\boxguard[20]
\begin{scenariobox}[Northgate Recreation Center]{cerulean}
Northgate keeps a record on every member. Here are eight of them.

\vspace{2pt}
\begin{center}
\small
\renewcommand{\arraystretch}{1.15}
\begin{tabular}{@{} l c l c c c @{}}
\toprule
\textbf{Member} & \textbf{Member ID} & \textbf{Membership} &
\textbf{Visits last month} & \textbf{Age} & \textbf{Group class?} \\
\midrule
Rosa  & 104 & Family  &  6 & 34 & Yes \\
Ellis & 118 & Student &  2 & 17 & No  \\
Wren  & 127 & Family  &  9 & 41 & Yes \\
Bo    & 133 & Senior  &  4 & 68 & No  \\
Tam   & 145 & Student & 14 & 19 & Yes \\
Odell & 152 & Family  &  1 & 52 & No  \\
Kira  & 166 & Student &  8 & 16 & Yes \\
Nash  & 179 & Senior  &  4 & 71 & Yes \\
\bottomrule
\end{tabular}
\end{center}
\end{scenariobox}

\vspace{0.08in}

\boxguard
\begin{notesbox}{Practice}
Every answer needs a reason --- ``it's words'' and ``it's numbers'' are not reasons.
\wordbank{8 \;\textbullet\; 5 \;\textbullet\; categorical \;\textbullet\; quantitative}
\begin{enumerate}[leftmargin=1.6em, itemsep=9pt, topsep=3pt]

  \item How many \textbf{individuals} does the table describe? \blank{1.8cm} \hfill
        How many \textbf{variables} does it record? \blank{1.8cm}
        \par\vspace{2pt}
        \emph{(``Member'' names the individuals --- it is not a variable being studied.)}

  \item Classify each variable, and say \textbf{how you know}.

        \begin{center}
        \renewcommand{\arraystretch}{1.5}
        \begin{tabularx}{\linewidth}{@{} p{3.2cm} p{3cm} X @{}}
        \toprule
        \textbf{Variable} & \textbf{Type} & \textbf{How I know} \\
        \midrule
        Member ID          & \blank{2.4cm} & \blank{6.2cm} \\
        Membership         & \blank{2.4cm} & \blank{6.2cm} \\
        Visits last month  & \blank{2.4cm} & \blank{6.2cm} \\
        Age                & \blank{2.4cm} & \blank{6.2cm} \\
        \bottomrule
        \end{tabularx}
        \end{center}

\end{enumerate}
\end{notesbox}

\boxguard[30]
\begin{notesbox}{Practice, continued}
\begin{enumerate}[leftmargin=1.6em, itemsep=6pt, topsep=2pt, start=3]

  \item Find the mean number of visits last month, then write it as \emph{data} ---
        one sentence naming the group and the units.

        \begin{work}
          \bar{x} &= \dfrac{6+2+9+4+14+1+8+4}{8} = \dfrac{48}{8} = 6 \text{ visits}
        \end{work}
        \par\writeline

  \item ``Group class?'' is categorical, so a mean is not available. Find the
        \textbf{percent} of these eight members who take a group class, then explain
        why a percent --- and not a mean --- is the right summary here.

        \begin{work}
          \dfrac{5}{8} &= 0.625 = 62.5\%
        \end{work}
        \par\writeline

  \item \textbf{Spiral --- Algebra 1.} The director asks, \emph{``How many of our $2000$ members
        use the pool?''} --- a statistical question, because the answer varies from member to
        member. A survey of $300$ members found that $45$ use the pool. Find the percent of
        those surveyed who use it, predict how many of all $2000$ do, and write the
        prediction as one sentence for the director.

        \begin{work}
          \dfrac{45}{300} &= 0.15 = 15\% \\
          0.15 \times 2000 &= 300 \text{ members}
        \end{work}
        \par\writeline

  \item A staff member computes the mean of the Member ID column --- correctly --- and gets
        $140.5$. She writes: \textit{``The average Northgate member is number $140.5$.''}
        What does that number actually tell you about Northgate's members? Use the type you
        gave Member ID in item~2 to explain.
        \par\writelines{2}

\end{enumerate}
\end{notesbox}

\boxguard[5]   % the two-line spiralbox needs about five baselines; a larger guard pushes it to a third page
\begin{spiralbox}
\textbf{Next lesson: The Data Cycle.} We name the four stages of the data cycle formally,
write questions a data set can actually answer, and organize collected answers into a
frequency table. Today's words --- \emph{individual}, \emph{variable}, \emph{categorical},
\emph{quantitative} --- are assumed from here on.
\end{spiralbox}

\end{document}
